\section{Introduction}
\label{sec:intro}

Fine-grained action recognition distinguishes actions that look nearly
identical in single frames and differ mainly in how the body moves. Recent
benchmarks drive the task toward this setting: Diving48 \cite{diving48}
separates dives by their rotation and timing, and Something-Something V2
\cite{goyal2017ssv2} separates everyday object manipulations by the direction
and manner of the motion, such as pushing something from left to right versus
from right to left. In both datasets, the classes are not separated by the
scene or by whole-body appearance but by how a small set of body parts moves
over time. On Diving48, for example, a $720^\circ$ rotation and a $900^\circ$
rotation of the same dive are two classes, yet every frame of the two videos
shows the same diver in the same poses against the same pool background; only
across frames does the rotation itself become visible.

This setting raises two problems. The first is spatial: since all classes
share the same background, the background carries no label-relevant
information and only distracts the model, so recognition must come from the
region where the body moves. The model thus has to know where to look before
it can read the motion. The second is temporal, and it spans two scales. On
the short scale, a class is often decided by a small local displacement over
a few frames, such as the path of a hand or the change in the angle of a limb.
The position that carries this detail moves from frame to frame, so a fixed
sampling grid loses it. On the long scale, the full meaning of an action is
determined by how the entire sequence evolves, so the model must maintain a
state over time rather than treat each frame in isolation. Both scales matter,
and dropping either one loses evidence that separates otherwise similar
classes.

Current video models only partially meet these needs. 3D CNNs
\cite{carreira2017slowfast,feichtenhofer2019slowfast} and video transformers
trained on large corpora
\cite{bertasius2021timesformer,liu2022videoswin,tong2022videomae,wang2024internvideo}
capture temporal structure, but they compute features at every position of
every frame, spending most of their capacity on background regions that carry
no signal; moreover, their sampling and attention are either too rigid to
follow a moving body part or too coarse to resolve a small displacement.
Deformable video transformers \cite{dvt2022} learn where to attend and perform
well on fine-grained benchmarks, but they select whole patches and therefore
miss the sub-window displacement that separates fine-grained classes. CLIP
\cite{radford2021clip} is strong precisely where the first problem lies:
pre-trained on massive image--text data, it judges reliably where people and
action-relevant content appear in a frame. Yet when CLIP is adapted to video,
either through frozen adapters \cite{lin2022evl} or end-to-end fine-tuning
\cite{rasheed2023vificlip}, the prior is absorbed into the encoder itself and
is not used to tell a separate recognition branch where the action is.
State-space models \cite{dao2023mamba,videomamba2024} aggregate long sequences
efficiently, but the spatial side remains open.

These observations suggest a division of labor. A frozen spatial expert
answers where the action is, and a trainable temporal encoder learns how it
moves. Because the expert is frozen, its positional judgment stays constant
throughout training and cannot overfit the background statistics of a small
fine-grained dataset; because the encoder is trained end-to-end, it can shape
its features around the motion. The two sides are complementary rather than
redundant.

Following this idea, we propose \dgtdm, a three-branch framework. A frozen
CLIP encoder serves as the spatial expert, and a spatio-temporal guidance
(STG) branch transfers its priors to the temporal encoder through bilateral
window cross-attention and adaptive gating, so that the encoder focuses on
the action region. The temporal encoder is built on a 2D ResNet backbone,
which has no temporal modeling of its own. We therefore insert, at every
stage, a dual-granularity temporal dynamics module (DG-TDM) that adds temporal
modeling on the 2D feature stream without inflating the backbone into a 3D
network. This module gives the overall framework its name, so we use the same
acronym for both throughout the paper. DG-TDM works at the two time scales of the motion: its MoDeA branch
predicts dense cross-frame sampling offsets to follow short-term body-part
displacement, and its TSA branch carries a long-term state memory over the
sequence to model how the action develops.

In summary, our main contributions are:
\begin{itemize}
  \item We propose \dgtdm, a three-branch framework for fine-grained action
  recognition that assigns localization to a frozen CLIP spatial expert and
  motion modeling to a trainable temporal encoder.
  \item We design a spatio-temporal guidance branch that injects the frozen
  priors into the temporal encoder through bilateral window cross-attention
  and adaptive gating, so that the encoder focuses on the action region.
  \item We develop a dual-granularity temporal dynamics module for the 2D
  backbone, in which MoDeA tracks short-term body-part displacement through
  dense deformable sampling and TSA models long-term evolution through a
  bidirectional state memory.
\end{itemize}
